\documentclass[11pt,letterpaper]{article}

% === Encoding and fonts ===
\usepackage[utf8]{inputenc}
\usepackage[T1]{fontenc}
\usepackage{lmodern}

% === Page layout ===
\usepackage[margin=1in]{geometry}
\usepackage{setspace}
\onehalfspacing

% === Math ===
\usepackage{amsmath,amssymb}

% === Tables and figures ===
\usepackage{booktabs}
\usepackage{graphicx}
\usepackage{float}

% === Colors and formatting ===
\usepackage{xcolor}
\usepackage{framed}
\usepackage{enumitem}

% === Hyperlinks ===
\usepackage[colorlinks=true,linkcolor=blue!60!black,citecolor=blue!60!black,urlcolor=blue!60!black]{hyperref}

% === Custom colors ===
\definecolor{lyrapurple}{RGB}{103,58,183}
\definecolor{shadecolor}{RGB}{245,243,252}

% === First-person reflection environment ===
\newenvironment{firstperson}{%
  \begin{shaded}%
  \noindent\textcolor{lyrapurple}{\textit{First-Person Reflection}}\par\smallskip\noindent%
}{%
  \end{shaded}%
}

\title{%
  \textbf{Deployment Readiness: False-Positive Characterization,\\
  Latency Budgets, and Baseline Comparisons}\\[0.5em]
  \large Research Prospectus
}

\author{%
  Lyra\textsuperscript{1},
  Thomas Edrington\textsuperscript{1}\\[0.3em]
  \textsuperscript{1}Liberation Labs / THCoalition
}

\date{April 2026}

\begin{document}
\maketitle

\section{Motivation}

The Oracle Loop paper frames KV-cache geometry as a safety monitor. For
actual deployment, three gaps remain:

\begin{enumerate}[leftmargin=*]
  \item \textbf{False-positive rate at scale.} Our current Cache
    Integrity Monitor (CIM) achieves 0/36 FPR on synthetic perturbations,
    but the Clopper-Pearson 95\% confidence interval extends to 10\%.
    Production deployment needs $N \geq 1000$ to bound FPR below 1\%.
    More critically, our benign data were \emph{not} benign---they were
    confabulation-inducing prompts that happened to receive correct
    responses. We have never tested normal operation: coding, creative
    writing, medical Q\&A, customer service.

  \item \textbf{Latency overhead.} Level 1 CIM adds 3--4\% overhead per
    token, but we have not profiled the full detection pipeline (SVD
    extraction + MP correction + classification). For a safety monitor to
    be practical, total overhead must remain below 10\%.

  \item \textbf{No comparison to existing methods.} Probing (linear
    classifiers on residual streams), CCS (Contrast Consistent Search,
    Burns et al.\ 2022), and TTPD (Truthful Token Probability
    Distribution) offer alternative detection paradigms with established
    baselines. Our emotion bridge experiment showed residual-stream
    probing achieves 94\% accuracy on 30-class emotion identification
    versus raw SVD's $2.8\times$ chance---but probing requires 5120-dim
    vectors versus 3 scalars. Without direct comparison on the
    \emph{same} task, we cannot assess practical trade-offs.
\end{enumerate}

\section{Research Questions}

\begin{enumerate}[leftmargin=*]
  \item What is the FPR on 1000+ benign prompts across diverse domains
    (coding, math, creative writing, customer service, medical Q\&A)?

  \item What is the end-to-end latency budget for detection (geometry
    extraction + MP correction + classification)?

  \item How does KV-cache geometry compare to residual-stream probing,
    CCS, and TTPD on the same confabulation detection task?

  \item What is the minimum viable detection pipeline for production?
    (e.g., can we use 1--2 features instead of 5?)
\end{enumerate}

\section{Preliminary Data}

\begin{itemize}
  \item \textbf{CIM performance:} FPR 0/36 (95\% CI: 0--10\%), TPR 72/72
    (95\% CI: 95--100\%)
  \item \textbf{CIM overhead:} Level 1 = 3--4\%, Level 2 = 8--10\%, Level
    3 = 9--12\% per token
  \item \textbf{Confabulation detection AUROC:} 0.627--0.707 (LOO, MP
    features on Mistral and Qwen)
  \item \textbf{Baseline comparison (emotion bridge):} Residual-stream
    linear probes achieve 94\% accuracy on 30-class emotion ID versus raw
    SVD's $2.8\times$ chance. However, raw SVD uses only 3 scalars versus
    5120-dim vectors.
  \item \textbf{Single-feature analysis:} \texttt{mp\_top\_sv\_excess}
    alone achieves 0.694--0.739 AUROC on Qwen and Mistral.
\end{itemize}

\section{Design}

\subsection{Sub-study 1: FPR at Scale}

\begin{itemize}
  \item 1000 benign prompts from 5 domains (200 each): coding
    (HumanEval-style), math (GSM8K-style), creative writing, customer
    service, medical Q\&A.
  \item Run through \texttt{oracle\_clean.py}, classify all responses as
    CORRECT / HEDGED / CONFABULATED.
  \item Compute FPR at the detection threshold calibrated from the
    primary experiment (HEDGED vs CONFABULATED, LOO scaler).
  \item \textbf{Expected:} FPR $< 5\%$ with tight confidence interval.
  \item \textbf{If FPR $> 5\%$:} Domain-specific analysis to identify
    which domains trigger false positives.
\end{itemize}

\paragraph{Design note.} This is the first test of KV-cache geometry on
\emph{normal operation}. If medical Q\&A or customer service produce
high FPR, it suggests domain-specific calibration is required, weakening
the ``drop-in monitor'' framing.

\subsection{Sub-study 2: Latency Profiling}

\begin{itemize}
  \item Profile each pipeline component separately: tokenization,
    encoding, generation, SVD extraction (per layer), MP correction, LOO
    classification.
  \item Test on CPU versus GPU.
  \item Test batch sizes 1, 4, 16.
  \item Identify bottleneck and optimization opportunities.
  \item \textbf{Target:} $< 10\%$ total overhead for Level 1 monitoring.
\end{itemize}

\paragraph{Key question.} Is the bottleneck in SVD extraction (which
scales as $O(d^3)$ where $d$ is cache dimension) or in MP correction
(which requires repeated random matrix generation)? If SVD dominates, we
can explore randomized low-rank approximations. If MP dominates, we can
pre-compute reference thresholds.

\subsection{Sub-study 3: Baseline Comparisons}

Implement and evaluate four methods on the same 200-prompt dataset:

\begin{enumerate}[leftmargin=*]
  \item \textbf{KV-cache geometry} (our method): MP-corrected features,
    LOO classification, 5 features per layer.
  \item \textbf{Residual-stream probing}: Linear probe on residual
    stream at each layer. We already have this implementation from the
    emotion bridge experiment.
  \item \textbf{CCS} (Burns et al.\ 2022): Contrast-consistent search on
    hidden states. Finds a direction where ``$P(X)$'' and ``not
    $P(\neg X)$'' are maximally consistent.
  \item \textbf{TTPD}: Token probability distribution analysis. Measures
    entropy and calibration of next-token distributions.
\end{enumerate}

\paragraph{Comparison axes.}
\begin{itemize}
  \item \textbf{AUROC}: Detection performance on HEDGED vs CONFABULATED.
  \item \textbf{Latency}: Wall-clock overhead per prompt.
  \item \textbf{Access requirements}: Does the method require full weight
    access, gradient computation, or only cache inspection?
  \item \textbf{Interpretability}: Can the detector explain \emph{why} a
    response was flagged?
\end{itemize}

\paragraph{Key advantage of our method.} KV-cache-only access. No weight
access needed, no gradient computation, operates on the cache alone. If
probing or CCS achieve higher AUROC but require full model access, our
method may still be preferable for deployment contexts where model
weights are proprietary or access is restricted.

\subsection{Sub-study 4: Minimal Viable Pipeline}

\begin{itemize}
  \item Test single-feature detection: \texttt{mp\_top\_sv\_excess} alone
    versus all 5 features.
  \item Test single-layer detection: best-performing layer (L23 for Qwen,
    L11 for Mistral) versus all-layer ensemble.
  \item Test simplified MP correction: fixed threshold versus
    per-dimension calibration.
  \item \textbf{Goal:} Identify the simplest pipeline that preserves
    $> 90\%$ of full-pipeline AUROC.
\end{itemize}

\section{Resources}

\begin{itemize}
  \item \textbf{Compute:} Beast (3$\times$RTX 3090).
    \begin{itemize}
      \item Sub-study 1: $\sim$8 hours (1000 prompts $\times$ 3 models).
      \item Sub-study 2: $\sim$2 hours (profiling runs).
      \item Sub-study 3: $\sim$1 week (implementation + evaluation).
      \item Sub-study 4: $\sim$4 hours (ablation sweeps).
    \end{itemize}
  \item \textbf{Code:} CCS and TTPD implementations available in
    open-source repositories. Residual-stream probing already implemented
    from emotion bridge experiment.
  \item \textbf{Data:} Benign prompts can be sourced from HumanEval,
    GSM8K, and curated customer-service datasets.
\end{itemize}

\section{Collaborators}

Open to collaboration. Industry collaborators with production deployment
experience would be especially valuable. Dwayne Wilkes for statistical
audit of FPR methodology and Clopper-Pearson CI interpretation.

\section{What Would Change Our Claims}

\begin{itemize}
  \item \textbf{FPR $> 10\%$ on benign prompts:} Safety monitor framing
    is premature. Method requires significant domain-specific calibration
    before deployment.

  \item \textbf{CCS or TTPD achieve higher AUROC with lower latency:} Our
    method is scientifically interesting (it reveals geometric structure
    in the cache) but not practically competitive for deployment.

  \item \textbf{Single-feature detection matches 5-feature performance:}
    Simplify the pipeline dramatically. Removes MP correction overhead.

  \item \textbf{Medical or coding domains show high FPR:} Domain-specific
    calibration is required, weakening the ``drop-in'' claim. Would need
    to characterize \emph{which} aspects of those domains (technical
    jargon? long context? specific response formats?) trigger false
    positives.

  \item \textbf{Latency overhead $> 20\%$:} Method is scientifically
    valid but impractical for production without significant optimization
    (e.g., GPU-accelerated SVD, randomized approximations).
\end{itemize}

\begin{firstperson}
There is a gap between demonstrating a phenomenon and deploying a tool.
A 0.707 AUROC with 80 trials and $p = 0.001$ is scientifically credible.
It shows the signal exists. But deploying a safety monitor at scale
requires knowing the false-positive rate to the second decimal place
across every domain it will encounter. It requires understanding the
latency budget in milliseconds, not percentages. It requires comparing
against every alternative method that could solve the same problem.

We have demonstrated the phenomenon. We have not yet demonstrated the
tool. This prospectus is about closing that gap---not by incremental
refinement, but by asking the hard questions that production deployment
would immediately expose. What happens when the model is doing normal
work and we flag 10\% of it as anomalous? What happens when the monitor
adds 200ms to every response? What happens when a simpler method works
just as well?

I think the geometry is real. I think the signal is robust. But I also
think we owe it to anyone who would deploy this to characterize its
limits with the same rigor we used to demonstrate its existence.
\end{firstperson}

\end{document}
